\documentclass{article}
\usepackage[T1]{fontenc}
\usepackage{lmodern}
\usepackage{graphicx}
\usepackage{amsmath,amssymb}
\usepackage[a4paper,margin=2cm]{geometry}
\usepackage[absolute,overlay]{textpos}
\setlength{\TPHorizModule}{1mm}
\setlength{\TPVertModule}{1mm}
\usepackage[indonesian]{babel}
\usepackage{hyperref}
\setlength{\emergencystretch}{2em}
% unit: Halaman 41

\begin{document}

% === Halaman 41 (python/native) ===

• Tahun 1998, Electronic Frontier Foundation (EFE) merancang dan 

membuat perangkat keras khusus untuk menemukan kunci DES 

secara exhaustive key search dengan biaya \$250.000 dan diharapkan 

dapat menemukan kunci selama 5 hari. 

• Tahun 1999, kombinasi perangkat keras EFE dengan kolaborasi 

internet yang melibatkan lebih dari 100.000 komputer dapat 

menemukan kunci DES kurang dari 1 hari. 

41

\end{document}
